\chapter*{부록: 재현 자료}\label{chap:appendix}
\addcontentsline{toc}{chapter}{부록: 재현 자료}

% Tables in the appendix are numbered A.1, A.2, ...
\setcounter{table}{0}
\renewcommand{\thetable}{A.\arabic{table}}

\ref{chap:experiments}장의 결과를 생성한 Jupyter 노트북과 실행 코드, 측정 데이터, 분석 결과는 다음 저장소에 있다.

\begin{center}
\url{https://github.com/yabsed/vlm-in-flash-jupyter}
\end{center}

Neuron Chunking의 greedy 방법, hot--cold 순서의 계산, SSD reader에는 원 논문의 공개 구현\footnote{\url{https://github.com/snuhcs/vlm-flash}}을 사용하였다. Projection별 행 예산의 탐색은 공개 구현에 없어 노트북에 구현하였다.

각 절의 결과를 생성한 노트북은 표~\ref{tab:notebooks}\와 같다. 경로는 이 저장소의 \texttt{jetson} 폴더를 기준으로 한다. \texttt{05\_01\_calibrate}는 노트북 컴퓨터에서 실행하였으며, 그 출력은 \texttt{results/05} 폴더의 같은 이름의 노트북에 있다.

\begin{table}[htbp]
    \centering
    \caption{절별 노트북}
    \label{tab:notebooks}
    \begin{tabular}{lll}
        \toprule
        절 & 노트북 & 내용 \\
        \midrule
        \ref{sec:exp-setup} & \texttt{01/01\_01\_measure\_chunk\_latency} & 읽기 비용 측정 \\
                            & \texttt{03/03\_01\_tune\_paper} & Paper의 후보 설정 \\
                            & \texttt{05/05\_01\_calibrate} & Hot--cold 순서와 행 예산 \\
                            & \texttt{05/05\_04\_report\_numbers} & 표~\ref{tab:projections}, 표~\ref{tab:selector-settings} \\
        \midrule
        \ref{sec:exp-measure} & \texttt{05/05\_02\_run\_vlm} & 모델 실행, 정확도와 프레임당 시간 측정 \\
        \midrule
        \ref{sec:exp-results} & \texttt{05/05\_03\_accuracy\_latency} & 그림~\ref{fig:accuracy-latency}, 표~\ref{tab:summary} \\
                              & \texttt{05/05\_04\_report\_numbers} & 표~\ref{tab:paired} \\
        \midrule
        \ref{sec:exp-paper} & \texttt{04/04\_01\_run\_vlm} & 준비 과정을 생략한 실행 \\
                            & \texttt{05/05\_03\_accuracy\_latency} & 표~\ref{tab:paper} \\
        \bottomrule
    \end{tabular}
\end{table}
